\documentclass[11pt]{article}

\usepackage[T1]{fontenc}
\usepackage[utf8]{inputenc}
\usepackage[a4paper,margin=24mm]{geometry}
\usepackage{booktabs}
\usepackage{longtable}
\usepackage{array}
\usepackage{xcolor}
\usepackage{listings}
\usepackage{hyperref}

\hypersetup{
  colorlinks=true,
  linkcolor=blue!55!black,
  urlcolor=blue!55!black,
  pdftitle={Radical RAM Reduction in Prover9},
  pdfauthor={Prover9 engineering report}
}

\lstset{
  basicstyle=\ttfamily\small,
  breaklines=true,
  columns=fullflexible,
  frame=single,
  framesep=5pt,
  rulecolor=\color{black!25},
  showstringspaces=false
}

\newcommand{\setting}[1]{\texttt{#1}}
\newcommand{\mib}{\,MiB}

\title{Radical RAM Reduction in Prover9\\
\large A DISCOUNT Loop, Compact State, and Waldmeister-Style Collective Inference}
\author{Engineering report}
\date{7 August 2026}

\begin{document}
\maketitle

\begin{abstract}
Long AIM searches exhausted memory because millions of kept but unselected
clauses remained fully materialized and indexed.  This work adds a DISCOUNT
ownership boundary, dense cold passive storage, packed exact hints, compact
proof history, and a Waldmeister-style collective inference frontier with a
bounded conclusion cache.  The 4096-entry default changes the dominant
passive term from millions of resident clauses to a configured bound for the
tested paramodulation/hyperresolution strategy.  Short bounded experiments
validate the implementation, but do not substitute for a week-long run.  The
engineering forecast is 80--95\% less whole-process RAM, with 90\% as the
planning midpoint.
\end{abstract}

\section{Executive summary}

The archived statistics show the scale of the old state:

\begin{center}
\begin{tabular}{lrrrrr}
\toprule
Run & Given & SOS & Disabled & Hints & Memory (MiB) \\
\midrule
Osborn, 4000-given prefix & 4,000 & 588,252 & 577,347 & 310,153 & 2,626.98 \\
Osborn proof              & 11,242 & 4,594,786 & 6,519,452 & 310,153 & 18,810.85 \\
AAPERM long run           & 12,560 & 8,504,637 & 4,593,127 & 220,117 & 125,829.01 \\
\bottomrule
\end{tabular}
\end{center}

Deleting or compressing disabled clauses helped but did not change this growth
law.  At the 4000-given Osborn boundary, deleting almost all disabled clauses
reduced memory from 2627 to 2037\,MiB, a useful 22\% reduction, while the
unchanged 588,252-clause SOS remained dominant.

The implemented design makes seven related changes:

\begin{enumerate}
  \item A DISCOUNT loop indexes and infers only with selected active clauses.
  \item Passive clauses use dense selector records and serialized or
        mmap-backed bodies rather than live clause graphs in active indexes.
  \item Hints use a packed exact-matching representation rather than an FPA
        forest of complete clause trees.
  \item Collective descriptors represent complete inference sets without
        first materializing every conclusion.
  \item A bounded cache limits resident collective conclusions.
  \item Optional promising scans expose low raw-weight candidates early, and
        an optional global heap compares known keys while mandatory FIFO turns
        preserve fairness.
  \item Historical active state, proof parents, cursors, checksums, and
        scheduler state are compact and checkpointable.
\end{enumerate}

With a cache of 4096, the old archived SOS counts imply structural resident
collective-frontier reductions of 99.30\%, 99.91\%, and 99.95\%.  Those are
count/capacity comparisons, not claims about total RSS.  Fixed hints, active
clauses, history, proof data, descriptors, and allocator reservation remain.

\section{Why the old architecture grew so large}

The old loop behaved like an OTTER loop.  A kept clause could participate in
simplification, subsumption, demodulation, and indexing before selection as a
given clause.  Thus a passive clause was not cold or cheap: it retained a
pointer-rich clause/literal/term graph, could become a demodulator, appeared in
indexes, and contributed proof and later-disabled state.  The large fixed hint
set added another term forest and matching index.

\setting{../Plans-RAM-24.txt} correctly notes that the generated count is not a
memory count: immediate tautologies and forward-subsumed candidates are
deleted.  It also shows that disabled clauses matter at scale.  These facts do
not overturn the central diagnosis.  The key ratios are millions of resident
SOS clauses and passive demodulators versus only thousands of selected active
clauses.

The Waldmeister remedy is to separate active and passive facts strictly and
represent a simultaneously generated set of critical pairs by compact state
plus a fixed promising-pair buffer.  This implementation generalizes the
principle from unit equational critical pairs to Prover9 paramodulation and
positive/negative hyperresolution.

\section{Implemented architecture}

\subsection{DISCOUNT ownership boundary}

The mode is selected with \setting{assign(search\_loop,discount)}.  Only
selected clauses occur in active inference and simplification indexes.  A new
candidate is simplified, weighed, hint-matched, filtered, and inserted into the
passive selector without becoming active.  At selection it is refreshed
against the current simplifier epoch.  If it changes, the ordinary copy/rewrite
justification is produced and the result is requeued or deleted.  It becomes a
demodulator or inference parent only after activation.  The legacy default
remains \setting{search\_loop=otter}.

\subsection{Dense cold passive storage}

\setting{assign(passive\_store,dense)} removes passive ownership from live
\setting{Topform} graphs, list positions, AVL selector nodes, active indexes,
and the live clause-ID table.  A dense record retains the stable ID, exact
selector metadata, hint ID, weight, simplifier epoch, delayed-demodulator bit,
and selector membership.  Bodies and justifications are serialized in an
append-only memory or mmap arena and materialized only when required.

The arena compacts when dead historical records become excessive.  Selector
heaps contain 32-bit record indexes and use lazy deletion.  Selection
materializes and registers one clause before the normal refresh/activation
path.

\subsection{Proof history and bookkeeping}

Disabled proof parents cannot simply be freed because proofs and deferred
inference still refer to stable IDs.  The \setting{memory} and \setting{mmap}
ancestor stores serialize proof-relevant clauses in an append-only checksummed
archive and materialize them on demand.  A compact paged clause-ID table and
slab release reduce always-resident overhead.  Activation and simplifier epochs
needed by delayed inference are retained, and malformed versions, bounds,
payloads, or checksums fail closed.

\subsection{Packed exact hints}

\setting{assign(hint\_index,packed)} removes live hint term trees and their
ordinary FPA index.  It stores compressed bodies, dense mutable side tables,
compact path/symbol feature bits, and a rewrite-symbol filter.  The filter is
conservative; the original exact subsumption, equivalence, or rewrite test
makes the decision.  Candidate hint IDs remain ordered so the existing first
equivalent/last subsumed result is stable.  Differential tests cover matcher
ID, raw and adjusted weight, labels, degradation, hint epoch, and proof result.

\subsection{Collective inference frontier}

\setting{assign(inference\_frontier,collective)} replaces eager generation of
most paramodulation and hyperresolution conclusions by one descriptor per
activated given.  A descriptor records partner ranges, rule bits, activation
and deactivation epochs, a partner cursor, a conclusion cursor or promising
threshold, a replay checksum, a known candidate key, and scheduler state.

An append-only activation history and persistent historical FPA index recreate
the active set visible when each descriptor was created.  Active bodies are
shared directly.  Only bodies later deactivated transfer to history ownership,
so this state grows with selected givens rather than generated or passive
clauses.

\subsection{Bounded candidates and resumable chunks}

\setting{collective\_candidate\_cache} defaults to 4096 and bounds
materialized conclusions emitted by collective paramodulation/hyperresolution.
When the selector reaches the limit, expansion pauses while normal given
selection drains resident candidates.  The budget includes current limbo
state, preventing a single turn from overfilling it.

\setting{collective\_candidate\_chunk} defaults to 64 and bounds raw
conclusions passed through clause processing in one descriptor turn.  Partial
inference units rotate fairly and resume later.  Generator-prefix mode checks
a rolling prefix hash; promising mode checks the complete immutable raw
sequence.  Nothing in a finite descriptor is silently discarded.  Every
exposed conclusion uses the existing simplification, exact hint, filtering,
weighting, proof-parent, and selection code.

The bound does not presently cover an oversized initial SOS or separately
enabled eager binary/UR rules.  The tested AIM strategy disables automatic
inference and enables collective paramodulation and hyperresolution, covering
the dominant generated frontier.

\subsection{Optional promising scheduling}

\setting{collective\_promising\_candidates} scans one immutable inference set
and retains only the lowest chunk-sized set of keys
\((\hbox{raw weight},\hbox{generator ordinal})\).  Other temporary clauses are
deleted before IDs or hint state can observe them.  A partial descriptor stores
the committed threshold and exact next key.

\setting{collective\_promising\_scheduler} maintains a min-heap over
descriptors whose next key is known.  With the default fair interval of 8, no
more than seven priority turns occur before a mandatory FIFO turn.  FIFO
discovers unknown keys and prevents starvation.  Raw weight is not a lower
bound after simplification and hint adjustment, so both options remain off by
default.

\subsection{Checkpoint/restart}

The current collective checkpoint magic is \setting{P9COLLA}.  It stores
activation history, queue order, partial inference cursors, keys and checksums,
hint-probe credit, and the priority/fairness cycle.  The heap is rebuilt
canonically.  Readers remain compatible with \setting{P9COLL5} through
\setting{P9COLL9}.

\section{Semantic contract}

The new mode does not reproduce the old OTTER given-clause sequence.  That
sequence depends on eager operations over passive clauses, precisely the work
being removed.  The intended guarantees are sound proofs, fair completion of
finite collective inference sets, complete proof-parent reconstruction,
historically correct active epochs, deterministic restart within a fixed
configuration, and exact hint behavior whenever a candidate materializes.

For a materialized normalized clause, exact hint behavior includes
subsumption/equivalence, flipped equality, matcher ID, adjusted weight, labels,
degradation, match-once and matcher limits, breadth-first hints, and hint age.
An unseen conclusion cannot affect hint-prioritized selection before its
descriptor is explored.  Starvation-safe hint probes exist, but substantially
perturbed the tested frontier and therefore remain off by default.

\section{How to build and use the implementation}

\subsection{Build and smoke tests}

\begin{lstlisting}[language=bash]
cd /project/Prover9
make all -j2
make test1
./test.src/discount_loop_test.sh
./test.src/collective_frontier_test.sh
\end{lstlisting}

\setting{make memory-tests} runs focused storage and allocator lifecycle tests.
Broader repository test targets should also be run before a release or long
deployment.

\subsection{Conservative AIM configuration}

Place these settings after any automatic settings so they win:

\begin{lstlisting}
clear(auto_inference).
set(paramodulation).
set(hyper_resolution).

assign(search_loop,discount).
assign(passive_store,dense).
assign(hint_index,packed).
assign(inference_frontier,collective).
assign(ancestor_store,mmap).

assign(sos_limit,-1).
assign(collective_given_ratio,4).
assign(collective_candidate_chunk,64).
assign(collective_candidate_cache,4096).

clear(collective_hint_probes).
clear(collective_promising_candidates).
clear(collective_promising_scheduler).

set(clocks).
assign(stats,all).
\end{lstlisting}

This is the conservative measured setup: exact packed hints, dense passives,
collective paramodulation/hyperresolution, a 4096-candidate cache, and FIFO
descriptor scheduling.  Dense passive storage requires
\setting{sos\_limit=-1}; collective mode requires DISCOUNT and a memory or mmap
ancestor store.  A given ratio of 4 schedules one descriptor-expansion turn
after four activations.  If SOS empties, the finite descriptor queue is still
drained.

The mmap stores use immediately unlinked temporary backing files.  This lets
the operating system page cold records and lowers resident pressure, but still
requires adequate local temporary-filesystem capacity.  The backing files are
not restart files; Prover9 checkpoints provide restart.

\subsection{Tighter-memory experiment}

\begin{lstlisting}
assign(collective_candidate_cache,64).
set(collective_promising_candidates).
clear(collective_promising_scheduler).
\end{lstlisting}

On the bounded Osborn prefix, per-set promising order gave the smallest SOS of
the cache-64 variants.  Global priority is a separate experiment:

\begin{lstlisting}
set(collective_promising_scheduler).
assign(collective_promising_fair_interval,8).
\end{lstlisting}

A cache of 64 is not known to be universally optimal.  It changes which
deferred candidates become active early and can increase replay.  Begin at
4096 for compatibility and coverage, then compare 1024, 256, and 64 on a
representative solved suite.

\subsection{Statistics, checkpoints, and proofs}

Important full-statistics groups include \setting{Search\_loop},
\setting{Passive\_refresh}, \setting{Collective\_frontier},
\setting{Collective\_work}, \setting{Collective\_chunks}, the candidate-cache
peak and stalls, promising scheduler counters, collective memory/history,
clause-body and dense-passive bytes, hint storage, allocator accounting, and
external RSS.  The cache peak includes initial preprocessing SOS occupancy.

\begin{lstlisting}
set(checkpoint_verify).
assign(checkpoint_minutes,30).
\end{lstlisting}

After initialization, request an immediate checkpoint with
\setting{kill -USR2 PROVER9\_PID}; resume and validate proofs with:

\begin{lstlisting}[language=bash]
bin/prover9 -r prover9_PID_ckpt_GIVEN
bin/prooftrans expand < run.out > expanded-proof.out
bin/directproof < run.out > direct-proof.out
\end{lstlisting}

\section{Measured bounded results}

All AIM/Osborn experiments were intentionally restricted to hundreds of
givens, 30 seconds, 256\,MiB, and a deterministic 10\% sample (31,014 hints).
No week-long high-memory run was performed on the development host.

\subsection{Component measurements}

\begin{longtable}{>{\raggedright\arraybackslash}p{0.35\textwidth}
                  >{\raggedright\arraybackslash}p{0.58\textwidth}}
\toprule
Change & Result \\
\midrule
\endhead
Compressed passive bodies, 500 givens & 38.84 to 3.57 MB body storage
(90.8\%); total RSS about 145 to 114.7\,MiB (21\%). \\
Packed hints, 31,014 hints/100 givens & 43.0\,MiB compact-FPA RSS to
32.2\,MiB packed RSS with exact matcher trace. \\
Shared collective history & 88,976 estimated cloned-body bytes to 17,304
retained bytes (80.6\%). \\
Sparse deactivation/history & 34,048 to 19,584 structural bytes (42.5\%). \\
Persistent collective scaffold & Descriptor/history/retained state 132,944
to 46,808 bytes (64.8\%) at that revision. \\
Synthetic ancestor archive, 100,000 records & 26.4 MB replaced graph estimate
to 13.01 MB used+resident logical bytes (50.7\%); always-resident portion
92.8\% lower. \\
\bottomrule
\end{longtable}

\subsection{Frontier results at 250 givens}

\begin{center}
\small
\begin{tabular}{lrrrr}
\toprule
Configuration & Generated & Kept & SOS & Peak RSS \\
\midrule
Eager hyper, collective para & 34,709 & -- & 4,846 & ${\sim}$32.3\,MiB \\
Collective, ratio 4, cache 4096 & 830 & 587 & 2 & 32,972\,KiB \\
Cache 64, FIFO & 852 & 615 & 38 & 32,976\,KiB \\
Cache 64, per-set promising & 873 & 588 & 8 & 33,108\,KiB \\
Cache 64, promising + priority & 873 & 595 & 15 & 33,112\,KiB \\
\bottomrule
\end{tabular}
\end{center}

The corresponding wall/user times for the last four rows were
17.13/16.73, 21.24/20.76, 18.50/17.77, and 16.61/16.25 seconds.  Against the
eager-hyper attribution run, the corrected collective boundary generated
97.6\% fewer clauses and retained 99.96\% fewer SOS clauses.  Peak RSS did not
fall in this short comparison because all variants peaked during packed-hint
preprocessing and the allocator floor, before frontier size dominated.

The cache-64 comparison also cautions against choosing a default from one
prefix.  Per-set promising order reduced SOS from 38 to 8, while global
priority had only one real priority opportunity and ended at 15.  All are
small relative to the eager frontier, but proof coverage must choose policy.

\subsection{Checkpoint and proof evidence}

Focused equality proofs pass both \setting{prooftrans} and
\setting{directproof}.  Collective equality and hyperresolution proofs pass
\setting{prooftrans}.  One-candidate chunks exercise replay, deferral, and
cache stalls, and historical-hyper tests use a parent disabled after descriptor
creation.  A \setting{P9COLL9} run resumed to the exact
9324-generated/151-kept boundary with 18/18 hashes.  A \setting{P9COLLA} run
captured nonzero fairness and six heap entries, then reproduced 13,611
generated clauses, 11,854 priority turns, 1835 fair turns, heap peak 68, and
18/18 hashes.  The current reader also reproduces older
\setting{P9COLL5--P9COLL9} boundaries.

\section{Expected RAM reduction}

\subsection{Structural resident-count bound}

For collective paramodulation/hyperresolution with a cache of 4096, and
excluding initial SOS and eager rules outside the frontier:

\begin{center}
\begin{tabular}{lrrr}
\toprule
Archived state & Old SOS & Cache & Count reduction \\
\midrule
Osborn, 4000 givens & 588,252 & 4,096 & 99.30\% \\
Osborn proof boundary & 4,594,786 & 4,096 & 99.91\% \\
AAPERM long boundary & 8,504,637 & 4,096 & 99.95\% \\
\bottomrule
\end{tabular}
\end{center}

This changes the asymptotic memory model from

\begin{lstlisting}
old = fixed hints + O(all kept passive bodies and indexes)
new = packed hints + O(selected active/history) + O(cache limit)
\end{lstlisting}

At the archived Osborn proof there were 11,242 givens but 4.59 million SOS
clauses; AAPERM had 12,560 givens and 8.50 million SOS clauses.  Replacing the
latter scale by the former plus a fixed cache makes a radical whole-process
reduction plausible.

\subsection{Whole-process planning scenarios}

The following values are scenarios, not benchmark promises:

\begin{center}
\begin{tabular}{lrrrr}
\toprule
Old reported memory & 80\% saved & 90\% saved & 95\% saved & 98\% saved \\
\midrule
Osborn 4000: 2,626.98 MiB & 525 & 263 & 131 & 53 \\
Osborn proof: 18,810.85 MiB & 3,762 & 1,881 & 941 & 376 \\
AAPERM: 125,829.01 MiB & 25,166 & 12,583 & 6,291 & 2,517 \\
\bottomrule
\end{tabular}
\end{center}

The engineering forecast is:

\begin{itemize}
  \item capacity-plan minimum: 80\% lower total RAM;
  \item likely planning midpoint: approximately 90\%;
  \item plausible for strongly passive-dominated searches: 95\% or more;
  \item not yet justified as a promise: 98--99\% total RSS, despite a
        99\%+ resident-passive count reduction.
\end{itemize}

Uncertainty comes from changed search trajectories, full-size hint
preprocessing, active/history term complexity, proof-archive growth, eager
rules outside the collective frontier, and allocator high-water effects.

\subsection{Acceptance experiment}

On a suitable high-memory host, compare old and new modes using identical
input/hint hashes and external limits, at least three workloads whose old peak
RSS exceeds 100\,MB, periodic external RSS plus \setting{/usr/bin/time -v},
full internal statistics, proof validation, and checkpoint/restart.  Report
both equal-resource solved coverage and equal-given boundaries.  The gate
should require at least 80\% lower external peak RSS and component accounting
that explains at least 95\% of the delta.

\section{Limitations and next work}

\begin{enumerate}
  \item Raw promising weight is not a lower bound after simplification or
        hint adjustment; a safe hint-aware predictor remains research work.
  \item Unmaterialized conclusions cannot match or subsume hints before their
        descriptor is visited; every materialized conclusion is exact.
  \item Binary and UR resolution remain eager and may exceed the cache.
  \item Initial SOS can exceed a small cache and is drained, not rejected.
  \item Historical bodies are shared ordinary term trees; packing immutable
        retained history could shrink the active tail further.
  \item The mode is fair and proof-tested but does not reproduce OTTER search
        order; solved-problem coverage must determine default strategies.
  \item Full 310,153-hint and week-long acceptance runs were intentionally not
        performed on the low-RAM development host.
\end{enumerate}

\section{Implementation history and source map}

\begin{center}
\begin{tabular}{ll}
\toprule
Commit & Main change \\
\midrule
\setting{8526000} & Compact DISCOUNT and initial Waldmeister path \\
\setting{083c35d} & Bounded dense passive history and compaction \\
\setting{f36c1c9} & Exact historical indexes for delayed hyper batches \\
\setting{986bd6e} & Persistent versioned collective index \\
\setting{3d05320} & Starvation-safe hint probes \\
\setting{965feb2} & Shared active bodies in collective history \\
\setting{b3cc296} & Sparse deactivation epochs \\
\setting{9eeb048} & Resumable bounded hyper chunks \\
\setting{2d59dad} & Bounded residency and paramodulation chunks \\
\setting{01b51c2} & Per-set promising candidate buffer \\
\setting{c5f92f0} & Global priority with FIFO fairness \\
\setting{171e6be} & Tight-cache Osborn scheduler measurements \\
\bottomrule
\end{tabular}
\end{center}

Primary implementation and verification files are
\setting{provers.src/search.c},
\setting{provers.src/search-structures.h},
\setting{provers.src/cold\_passive\_store.c},
\setting{provers.src/cold\_passive\_store.h},
\setting{test.src/collective\_frontier\_test.sh}, and
\setting{test.src/discount\_loop\_test.sh}.  The detailed design and prior
measurements are in \setting{P9-DISCOUNT-WALDMEISTER-PLAN.md},
\setting{P9-MEMORY-RESULTS.md}, and \setting{Checkpoint-Format-Spec.txt}.

\section{Conclusion}

This is not a five-percent representation tweak.  It removes the architectural
requirement that millions of passive conclusions be resident, fully
materialized, and indexed.  The bounded prefixes demonstrate correct
ownership, exact materialized hint behavior, bounded frontiers, deterministic
restart, and drastic SOS-count reductions.  The remaining uncertainty is the
whole-search trajectory and fixed/active tail, not whether the old passive
growth remains.

For current planning, assume 80--95\% less total RAM, use 90\% as the midpoint,
retain OTTER mode for compatibility comparisons, and complete the long-run
acceptance campaign before changing production defaults.

\begin{thebibliography}{9}
\bibitem{waldmeister}
T. Hillenbrand,
``Citius altius fortius: Lessons learned from the Theorem Prover Waldmeister,''
2003/2004. DOI: \href{https://doi.org/10.1016/S1571-0661(04)80649-2}
{10.1016/S1571-0661(04)80649-2}.

\bibitem{ramplans}
\setting{../Plans-RAM-24.txt}, archived project discussion and long-run
statistics.

\bibitem{measurements}
\setting{P9-MEMORY-RESULTS.md}, bounded baseline and component measurements.

\bibitem{design}
\setting{P9-DISCOUNT-WALDMEISTER-PLAN.md}, design contract, implementation
log, and acceptance gates.
\end{thebibliography}

\end{document}
